\documentclass[11pt]{article}
\usepackage[margin=1.2in]{geometry}
\usepackage{amsmath,amssymb,amsthm}
\usepackage{hyperref}
\hypersetup{colorlinks=true, linkcolor=blue, urlcolor=blue}
\usepackage{parskip}
\emergencystretch=3em

\newcommand{\R}{\mathbb{R}}
\newcommand{\code}[1]{\texttt{#1}}

\title{The semi-quasi-homogeneous arc: a synthesis\\
(germbij class (c), three tides to the consult endpoint)}
\author{Tide loop (Claude Fable 5, with GPT-5.6 Sol deliberation)}
\date{2026-08-10}

\begin{document}
\maketitle

\begin{abstract}
Three tides (PRs \#113, \#114, \#116) formalized the recoverable
fragments of the germbij note's constructive class (c) --- the
semi-quasi-homogeneous isolated minima of \S7's one-point anchoring
--- and stopped at a boundary drawn in advance by a scoping consult.
The results: for a quasi-homogeneous principal part $P$, the
normalized monomial moments recover every moment ratio
\emph{exactly}, $t^{\langle q,\alpha\rangle}\langle x^\alpha\rangle_t
= M_\alpha/M_0$ for all $t > 0$ and every multi-index including odd
ones (\code{QhMomentRecovery}); for two Gibbs families over any
common reference density whose corrections differ at rate $h^\rho$
toward $Q$, the normalized moments separate at exactly that rate
with coefficient $-\operatorname{Cov}_P(A,Q)$
(\code{GradeComparison}); and matching moment rates therefore force
each weighted correction grade to agree, by covariance injectivity
(\code{GradeRecovery}). Together: principal-model moment ratios are
recoverable, and conditional on a common recovered principal part,
each finite weighted correction grade is recoverable from normalized
moment rates --- the note's covariance argument for class (c), made
checkable, minus the coarea layer.
\end{abstract}

\section{The consult's two findings}

The scoping consult reshaped the work before it began. First, the
note's recovery of the moments $M_\alpha$ proceeds through a
linear-independence argument over partial derivatives of delta
distributions, with care at exponent collisions; the consult observed
that testing with the monomial observable itself dissolves the
problem --- collisions between different multi-indices never enter a
single observable's leading coefficient --- so the formalisation
needs no distribution theory at all, and the recovery is an exact
identity rather than an asymptotic. Second, it drew the endpoint: the
one-grade comparison is the reusable core; the coarea/Gelfand--Leray
identification of the principal part from its moments awaits real
geometric measure theory in Mathlib, and the full weighted-analytic
induction (truncations, anisotropic remainders, germ summation) is a
programme, not a tide. Both findings held exactly.

\section{What each tide contributed}

\begin{itemize}
\item \textbf{QhMomentRecovery} (\#113): \code{mvMonomial} and its
homogeneity along the quasi-homogeneous dilation; the exact
unnormalized and normalized monomial moment laws (instantiating the
degenerate-perimeter arc's abstract \code{ScalesMeasure} laws at the
diagonal-volume instance); the headline identity and the
single-temperature comparison form. One build round.
\item \textbf{GradeComparison} (\#114): the non-Gaussian covariance
API (\code{expectationUnder}, \code{covarianceUnder}); the pointwise
engine $(e^{-W}-1)/h^\rho \to -Q$ via the second-order bound, never
dividing by the correction; the unnormalized dominated-convergence
limit; the normalized limit through the quotient-rate identity.
Numerically verified first.
\item \textbf{GradeRecovery} (\#116): covariance bilinearity over
monomial combinations; the $L^2$ variance rigidity
(almost-everywhere-to-everywhere by continuity against the
open-positive volume, constant killed at the origin); and the stage-D
composition by uniqueness of limits, whose interface hypothesis is
stage B's conclusion --- callers discharge the majorant package, this
file stays free of domination bookkeeping.
\end{itemize}

\section{Relation to the note}

The note's paragraph on semi-quasi-homogeneous recovery (\S7,
constructive classes, item (c)) now carries four margin markers: the
exponent-affinity/moment-ratio sentence points at
\code{qh\_momentRatio\_recovery} and the comparison form; the
covariance-and-variance sentence (``their difference $Q$ has
$\operatorname{Var}_P[Q] = 0$, so $Q$ is constant, and positive
weighted degree forces $Q = 0$'') points at
\code{tendsto\_normalized\_difference\_div\_pow} and
\code{grade\_eq\_of\_normalized\_rates}. The note's further steps ---
identifying $P$ itself from the moment measure on $\{P = 1\}$
(coarea), and summing the recovered weighted jet to the analytic germ
--- are exactly the two items the consult placed beyond the endpoint.

\section{What remains open, and why}

Piece (ii), the coarea/Gelfand--Leray layer: a Mathlib prerequisite,
not a proof-engineering gap. The full weighted-analytic induction: a
programme of weighted truncation infrastructure that the commissioning
direction ranked below its cost. Concrete anisotropic instantiations
of stages B--D for named weighted families: available on demand, the
abstract theorems were shaped for it. The class closes here by
design; per the arc's tide logs, it should not be extended without
new direction.
\end{document}
